\subsubsection{Norme du commutateur des projections arithmétiques}

On obtient le préfacteur électronique sans utiliser \(\alpha\), la torsion
ou une unité dimensionnelle. Le domaine est la troisième réalisation de APP,
\(\mathcal A_3=\{1,\ldots,9\}^3\), munie de la mesure uniforme.
L’espace de fonctions est \(\ell^2(\mathcal A_3;\mathbb R)\), muni du produit
scalaire
\[
 \langle f,g\rangle=\frac1{729}\sum_{x\in\mathcal A_3}f(x)g(x).
\]
Dans ce bloc, \([P,Q]=PQ-QP\) désigne le commutateur additif d’opérateurs.
Nous définissons
\[
 \Sigma(x)=\rho_9(x_1+x_2+x_3),\qquad
 \Pi(x)=\rho_9(x_1x_2x_3).
\]
Les espérances conditionnelles \(P_\Sigma\) et \(P_\Pi\) sont les
projections orthogonales sur les fonctions constantes sur chaque fibre
additive et multiplicative, respectivement. Ainsi, pour une fonction \(f\),
\[
 (P_\Sigma f)(x)
 =\frac{1}{|\Sigma^{-1}(\Sigma(x))|}
   \sum_{y:\,\Sigma(y)=\Sigma(x)}f(y),
\]
et de même pour \(P_\Pi\). Les deux projections agissent sur le même espace ; elles ne représentent pas deux supports arithmétiques indépendants.

\begin{lemma}[Deux projections et angles principaux]
Si \(s\in[0,1]\) est une valeur singulière de l'appariement entre des bases
orthonormales des images de deux projections \(P,Q\), son bloc non
trivial contribue à \(\|[P,Q]\|\) par
\(s\sqrt{1-s^2}\).
\end{lemma}
\begin{proof}
Dans le plan principal correspondant, on peut écrire
\[
 P=\begin{pmatrix}1&0\\0&0\end{pmatrix},\qquad
 Q=\begin{pmatrix}
 s^2&s\sqrt{1-s^2}\\
 s\sqrt{1-s^2}&1-s^2
 \end{pmatrix}.
\]
Leur commutateur est
\[
 [P,Q]=s\sqrt{1-s^2}
 \begin{pmatrix}0&1\\-1&0\end{pmatrix}.
\]
La matrice finale est orthogonale, donc sa norme est un. Les blocs communs ou orthogonaux ont un commutateur nul. La décomposition orthogonale en plans principaux donne l'affirmation et réduit la norme totale au maximum de ces contributions.
\end{proof}

\begin{proposition}[Norme exacte d'incompatibilité]
Sur \(\ell^2(\mathcal A_3)\),
\begin{equation}
 \|[P_\Sigma,P_\Pi]\|=\frac{\sqrt3}{4}.
 \label{ivloc:norma}
\end{equation}
\end{proposition}
\begin{proof}
Soit \(N_{ab}=|\Sigma^{-1}(a)\cap\Pi^{-1}(b)|\). L'énumération entière
des trois symboles fournit
\[
 N=9\begin{pmatrix}
 0&1&1&0&1&1&0&1&4\\
 1&0&1&1&0&1&1&0&4\\
 1&0&2&0&1&2&0&0&3\\
 0&1&1&0&1&1&0&1&4\\
 1&0&1&1&0&1&1&0&4\\
 0&0&2&1&0&2&0&1&3\\
 0&1&1&0&1&1&0&1&4\\
 1&0&1&1&0&1&1&0&4\\
 0&1&2&0&0&2&1&0&3
 \end{pmatrix}.
\]
Chaque ligne a pour somme \(81\) ; les sommes des colonnes sont
\[
 (m_1,\ldots,m_9)=(36,36,108,36,36,108,36,36,297).
\]
Les indicatrices normalisées des fibres ont pour matrice
d'appariement \(C_{ab}=N_{ab}/\sqrt{81m_b}\). Par conséquent,
\[
 M=CC^{\mathsf T},\qquad
 M_{ac}=\sum_{b=1}^9\frac{N_{ab}N_{cb}}{81m_b}
\]
est une matrice rationnelle dont le spectre se compose des carrés des
valeurs singulières requises.

La multiplication exacte montre que \((1,\ldots,1)\),
\((1,-1,0)^3\) et \((1,1,-2)^3\) sont des vecteurs propres de valeurs propres
\(1\), \(1/4\) et \(7/132\), respectivement. Le sous-espace des
vecteurs à support dans \(\{3,6,9\}\) dont la somme est nulle a pour dimension
deux et pour valeur propre \(1/18\). Les vecteurs de somme nulle à support dans
\(\{1,4,7\}\) ou dans \(\{2,5,8\}\) forment un noyau de dimension
quatre. Ces sous-espaces sont indépendants et la somme de leurs dimensions vaut
neuf. Le spectre complet est donc démontré
\[
 \operatorname{spec}(M)=
 \left\{1,\frac14,\frac7{132},\frac1{18},\frac1{18},0,0,0,0\right\}.
\]
Le maximum de \(\lambda(1-\lambda)\) sur cet ensemble est \(3/16\), atteint en \(\lambda=1/4\). Le lemme précédent donne \(\sqrt{3/16}=\sqrt3/4\).
\end{proof}

Le mode \((1,-1,0)^3\) est précisément la distribution des valeurs du
sélecteur tritique de phase dans la coordinatisation nonadique. La norme scalaire retient
une mesure d'incompatibilité, mais ne conserve pas à elle seule
l'orientation de ce mode et ne reconstruit pas les deux projecteurs. Cette mémoire
demeure dans l'état de route qui accompagne la lecture scalaire.

\begin{proposition}[Algèbre du bloc électronique local]
\label{ivloc:algebra}
Dans le plan principal qui atteint \eqref{ivloc:norma}, soient
\[
 P=\begin{pmatrix}1&0\\0&0\end{pmatrix},\qquad
 Q=\begin{pmatrix}1/4&\sqrt3/4\\\sqrt3/4&3/4\end{pmatrix},\qquad
 S=2P-I,\qquad J=\frac4{\sqrt3}[P,Q].
\]
Alors \(S^2=I\), \(J^2=-I\), \(SJ=-JS\), et l'algèbre réelle engendrée par \(S,J\) est \(M_2(\mathbb R)\), réalisation de \(\mathrm{Cl}_{1,1}\). En outre,
\[
 n_\pm=\frac{S\pm J}{2},\qquad
 n_\pm^2=0,\qquad n_+n_-+n_-n_+=I.
\]
\end{proposition}
\begin{proof}
On a \(S=\operatorname{diag}(1,-1)\) et
\(J=\bigl(\begin{smallmatrix}0&1\\-1&0\end{smallmatrix}\bigr)\).
La multiplication vérifie les trois relations. Les matrices \(I,S,J,SJ\) sont linéairement indépendantes et forment une base de \(M_2(\mathbb R)\), ce qui prouve l'affirmation sur l'algèbre. Enfin, \((S\pm J)^2=S^2+J^2\pm(SJ+JS)=0\), et la somme des produits croisés est \((S^2-J^2)/2=I\).
\end{proof}

Dans ce bloc coexistent une direction elliptique, une direction hyperbolique et deux
directions nulles paraboliques. Le résultat décrit une représentation
matricielle locale de l'arithmétique des deux feuillets. Il n'identifie pas toute APP
à \(M_2(\mathbb R)\), ni à un carré latin quantique ; de telles
affirmations exigeraient d'autres domaines et d'autres relations. L'exponentielle
n'appartient pas non plus exclusivement au régime parabolique : elle peut
s'appliquer à n'importe lequel des générateurs de ce bloc.
